Die Relevanz der Arbeit ergibt sich aus gesellschaftlichen, rechtlichen und technologischen Aspekten. Gesellschaftlich ist der Zugang zu verständlichen Informationen eine Vorraussetzung für Teilhabe. Rechtlich ist barrierefreie Kommunikation unter anderem durch die UN-Behindertenrechtskonvention sowie medienrechtliche Vorgaben verankert \cite{}.

Technisch leistet die Arbeit einen Beitrag zur Weiterentwicklung automatisierter Textvereinfachung, indem sie quantitative Metriken, regelbasierte Validierung und LLM-Feedbackschleifen kombiniert. Insbesondere im Bereich der Leichten Sprache Plus existieren bislang nur wenige systematische Ansätze zur automatisierten Qualitätssicherung.